\section{结论}
\label{sec:conclusion}

扩大监测规模的收益取决于估计目标。估计当期再生数时，估计量 $C_t/(\rho n)$ 在每个有限规模上都达到 Cramér--Rao 下界，规模越大越准。估计跨期平均水平时，方差为 $v_R+A/m$，环境方差 $v_R$ 构成一个规模无法突破的底板；交叉点 $\mx=A/v_R$ 处两项各占一半，此后规模翻倍只能把方差降低不到四分之一。

这一结构给出了三条设计含义。扩大覆盖面比提高报告率更能压低方差，但对住院数据两者差别很小。多期预算下，最优单期规模为 $\sqrt{(c_0/c_1)\mx}$。小规模时判断疫情是否增长，应使用精确检验，并报告“此后始终保持目标功效”的规模。

美国流感住院数据支持“常数加规模反比”的误差形状。流行初期的经验交叉点为 22.6 例/周，与情景区间同一量级，并在多种替代规格和两组外部数据中保持规模项为正。但截距与斜率都在情景隐含值的三倍以上，部分来自分母噪声，部分来自环境的时间相关。经验交叉点因此只能作为量级参考，不能当作结构参数的估计。

对监测实践而言，预警需要尽可能大的病例规模。趋势估计的规模到了交叉点量级以后，增加观测期数或汇总相邻辖区更有效。流行阶段转换时，应当先更新基准模型。
